\chapter{Практические задачи}
\label{ch:labs}

Практический блок курса охватывает проверку гипотез, множественные сравнения, регрессию, временные ряды, причинность и марковские модели.
Формулировки задач приведены в настоящей главе.
Численные таблицы и описание переменных --- в приложении~\ref{app:lab-data}.
Тексты, изображения и аудиозаписи, не помещающиеся в приложение, указаны сносками к соответствующим задачам.
Подмножество задач для сдачи назначает преподаватель (вариант выдаётся скриптом в репозитории курса, \PSADRepo{}).

Каждое задание оформляется в электронном рабочем тетради курса: в начале --- номер и формулировка, далее --- графики, расчёты и выводы. Язык вычислений не фиксируется.
Разобранные шаги семинаров (загрузка данных, тесты, модели) --- в \S\S\ref{sec:practice-ch2}--\ref{sec:practice-ch8}; здесь --- только постановки лабораторных работ.

\begin{table}[htbp]
\centering
\caption{Соответствие лабораторных работ, глав и раздела «Практика и семинар»}
\label{tab:lab-theory-practice}
\footnotesize
\begin{tabularx}{\linewidth}{|l|X|X|}
\hline
\textbf{Блок} & \textbf{Главы теории} & \textbf{Разбор в книге} \\
\hline
Лаб.\ 1, зад.\ 1 (GOF, мощность) & \ref{ch:hypothesis}, \S\ref{sec:gof-ch2} & \ref{sec:practice-ch2} \\
Лаб.\ 1, зад.\ 1.2--1.4 & \ref{ch:hypothesis}, \ref{ch:multiple} & \ref{sec:practice-ch2}, \ref{sec:practice-ch3} \\
Лаб.\ 1, зад.\ 2 (регрессия) & \ref{ch:regression} & \ref{sec:practice-ch5} \\
Лаб.\ 2, ряды и SPRT & \ref{ch:timeseries} & \ref{sec:practice-ch6} \\
Лаб.\ 2, DAG и HMM & \ref{ch:causality} & \ref{sec:practice-ch7} \\
Лаб.\ 2, байес (при необходимости) & \ref{ch:bayes} & \ref{sec:practice-ch8} \\
\hline
\end{tabularx}
\end{table}

\section{Лабораторная работа~1}
\label{sec:lab1-tasks}

Первый блок опирается на главы~\ref{ch:hypothesis}--\ref{ch:regression}.
Теоретические задачи требуют полного аналитического вывода: все переходы описаны, выкладки и преобразования явно посчитаны (см.~\S\ref{sec:report-requirements}).

\subsection{Задание 1}


\subsubsection{Задача 1.1}

Проверить мощность и консервативность критериев Лиллиефорса, Харке--Бера, Шапиро--Уилка для выборок из следующих распределений:
\begin{itemize}
    \item Нормальное
    \item Лапласа
    \item Стьюдента
    \item Усеченное нормальное распределение (модуль каждого элемента выборки не превосходит 2)
\end{itemize}

\subsubsection{Задача 1.2}

Задана выборка описаний 150 экземпляров ириса разных видов \citep{fisher1936}. Описание каждого ириса состоит из четырёх признаков:
\begin{itemize}
    \item длина наружной доли околоцветника (sepal length)
    \item ширина наружной доли околоцветника (sepal width)
    \item длина внутренней доли околоцветника (petal length)
    \item ширина внутренней доли околоцветника (petal width)
\end{itemize}
Требуется определить, насколько в среднем различается каждая из этих характеристик между разными видами.
Для каждой характеристики выбрать подходящий размер эффекта \citep{ellis2010} и посчитать соответствующее значение.

\subsubsection{Задача 1.3}

Проверить мощность и консервативность критерия Уилкоксона о равенстве медиан для выборок вида:
X1: ~ $\alpha$ * N(0,1) + (1-$\alpha$) * N(2, 4);

X2: ~ $\alpha$ * N(0,1) + (1-$\alpha$) * N(2, 4) + $\delta$.

Здесь $\delta$ --- сдвиг, дающий возможность разделить выборки X1 и X2.
Изучить зависимость от $\alpha$ и $\delta$.
Распределение является гауссовой смесью, а не суммой независимых гауссовых величин. Для генерации выборки достаточно на каждом шаге выбирать компоненту смеси с вероятностями $(\alpha, 1-\alpha)$ и сэмплировать из соответствующего нормального закона.

\subsubsection{Задача 1.4}

Частота слов в тексте часто описывается законом Ципфа \citep{zipf1949}. В выборке из $N$ различных слов вероятность слова ранга $k$ аппроксимируется $P(k) \propto k^{-a}$, где $a>0$ --- параметр распределения.
Проверить формальным критерием, насколько закон Ципфа описывает частоты слов в двух корпусах: романе «Анна Каренина» и новостном корпусе OPUS-WMT-News.\footnote{Тексты корпусов: \url{http://az.lib.ru/t/tolstoj_lew_nikolaewich/text_0080.shtml} и \url{https://object.pouta.csc.fi/OPUS-WMT-News/v2019/mono/ru.txt.gz}.}
Можно ли утверждать, что параметры $a$ для этих двух корпусов совпадают?

\subsection{Задание 2}


\subsubsection{Задача 2.1}

Пусть задано два вектора ответов: $y$ --- истинный вектор ответов для некоторой выборки, а также есть вектор ответов $\hat{y}$ --- некоторой предсказательной модели. Наблюдатель хочет проверить гипотезу о том, что ровно в 25\% случаев модель даёт заниженные оценки.
Предложите метод проверки данной гипотезы: запишите задачу формально,
предложите статистику для решения данной задачи на уровне значимости $\alpha=0{,}05$.
Также найдите зависимость мощности данного критерия в зависимости от истинного процента заниженных ответов.
Все выкладки должны быть сделаны аналитически, без использования компьютера (допускается использование компьютера для подстановки численных значений в финальную формулу).

\subsubsection{Задача 2.2}

Рассмотрим фирму, которая занимается продажей лотерейных билетов. Правила лотереи следующие: все билеты являются выигрышными с вероятностью $p$. Билеты продаются до тех пор, пока хоть один человек не выиграет (гарантируется, что как только билет купили и он выигрышный, больше билетов не продают).
Из-за вспышки коронавируса все заводы закрылись, а с ними и знание заветного $p$ также пропало. Фирма хочет восстановить $p$, имея отчёт о продажах билетов за последние 5~месяцев: 8, 12, 7, 6 и 12~билетов.
Требуется:
\begin{enumerate}
    \item Оцените методом максимального правдоподобия параметр $p_0$.
    \item Записать задачу формально для проверки гипотезы $p=p_0$ по эксперименту с $n=100$ вскрытыми билетами, повторённому 10~раз (результаты: 13, 8, 11, 10, 11, 12, 7, 9, 10, 9).
    \item Предложить статистику для решения данной задачи.
    \item Получить приближённо нулевое распределение данной статистики.
    \item Записать явно правило принятия решения для обеспечения уровня значимости $\alpha=0{,}05$.
    \item Проверить гипотезу по записанному критерию для данных из условия. Противоречат ли они гипотезе?
\end{enumerate}
Все выкладки должны быть сделаны аналитически, без использования компьютера (допускается использование компьютера для подстановки численных значений в финальную формулу).

\subsubsection{Задача 2.3}

Известно, что электричка "Вашингтон-Петушки" аварийно останавливается раз в несколько дней.
Аналитики РЖД проанализировали, сколько дней электричка едет без поломок, и составили выборку:
x = (30, 52, 43, 36, 48, 35, 36, 40, 37, 45).
РЖД хочет проверить гипотезу, что дисперсия распределения равна 9 против правосторонней альтернативы.
Требуется:
\begin{enumerate}
    \item Ввести предположение, каким распределением описывается данная выборка.
    \item Записать задачу формально.
    \item Предложить критерий для оценки дисперсии распределения.
    \item Проверить гипотезу о значении дисперсии распределения для уровня значимости $\alpha=0{,}05$ аналитически.
    \item Вывести и получить доверительный интервал для значения дисперсии при $\alpha=0{,}05$.
\end{enumerate}
Все выкладки должны быть сделаны аналитически, без использования компьютера. (допускается использование компьютера для подстановки численных значений в финальную формулу)

\subsubsection{Задача 2.4}

Одеяла с электрообогревом применяются в хирургии для восстановления температуры тела пациента после операции. 
Имеются два вида одеял: стандартный (b0) и экспериментальный (b1). 
Для 14 пациентов известно время, за которое нормальная температура тела восстанавливается при использовании одеяла каждого из видов.
Как понять, отличаются ли экспериментальные одеяла от стандартного?
Требуется:
\begin{enumerate}
    \item Записать задачу формально в виде проверяемой гипотезы и альтернативы.
    \item Предложить не менее двух критериев и соответствующих статистик для проверки этой гипотезы и описать, при каких дополнительных условиях стоит применять тот или иной критерий, а также преимущества и недостатки каждого критерия.
    \item Аналитически выразить достигаемый уровень значимости каждого критерия на выборке или описать, как его получить с помощью табличных данных.
\end{enumerate}
Все выкладки должны быть сделаны аналитически, без использования компьютера (допускается использование компьютера для подстановки численных значений в финальную формулу).

\subsection{Задание 3}


\subsubsection{Задача 3.1}

Изучить поведение FDR для модельного эксперимента из гл.~\ref{ch:multiple} (табл.~\ref{tab:model-mht}, $m=200$, $m_0=150$). Рассмотреть случаи, когда число гипотез $m$ варьируется от 200 до 100000, для следующих поправок:
\begin{itemize}
    \item Без поправок
    \item Метод Холма
    \item Метод Бенджамини--Хохберга
\end{itemize}

\subsubsection{Задача 3.2}

Для каждого игрока NBA известны (данные --- приложение~\ref{app:data-nba}):
\begin{itemize}
    \item количество успешных бросков в домашних играх (\texttt{score\_home})
    \item количество бросков в домашних играх (\texttt{atm\_home})
    \item количество успешных бросков в гостевых играх (\texttt{score\_away})
    \item количество бросков в гостевых играх (\texttt{atm\_away})
\end{itemize}
Требуется определить, есть ли разница в успехе бросков у игроков в домашних и гостевых играх.
У какого процента игроков разница в успехе существенна?

\subsubsection{Задача 3.3}

На наборе Wine \citep{forina1986} предложить метод выбора наиболее важных признаков для логистической регрессии на основе изученных методов прикладной статистики.
Осуществить выбор.

\subsubsection{Задача 3.4}

Даны результаты работы двух машинных переводчиков на небольших выборках переводов для разных языковых пар.\footnote{Исходные тексты: \PSADPath{labs/lab1/data/mt}.}
Стандартная оценка качества перевода --- метрика BLEU \citep{papineni2002} (реализация в стандартных пакетах NLP).
Требуется определить:
\begin{itemize}
    \item превосходит ли один переводчик в среднем по парам второй переводчик по переводу
    \item связано ли качество перевода для разных языковых пар у двух переводчиков?
\end{itemize}
При подсчете BLEU учитывать только слова, регистр не учитывать.

Название файлов имеет формат \texttt{lang1\_lang2\_<translator\_id>.txt}.
\texttt{lang\_1}, \texttt{lang\_2} --- языки (перевод с \texttt{lang\_1} на \texttt{lang\_2}).
\texttt{gold} --- эталонный вариант, с которым сравнивается перевод от систем машинного перевода.

\subsection{Задание 4}


\subsubsection{Задача 4.1}

Рассмотрим данные по числу заболевших и выздоровевших от коронавируса в разных странах (приложение~\ref{app:data-corona}). Требуется проверить гипотезу о тому, что число выздоровевших
людей в странах не зависит от числа заболевших в стране. 
Требуется:
\begin{enumerate}
    \item записать задачу формально;
    \item предложить статистику для решения данной задачи;
    \item записать приближенно нулевое распределение данной статистики;
    \item записать явно правило принятия решения на основе статистики и нулевого распределения для обеспечения уровня значимости $\alpha=0{,}05$;
    \item проверить гипотезу по записанному критерию, для данных из условия. Противоречат ли они гипотезу?
    \item на уровне значимости $\alpha=0{,}05$ найти зависимость мощности критерия в зависимости от истинного значения статистики.
\end{enumerate}
Все выкладки должны быть сделаны аналитически, без использования компьютера. (допускается использование компьютера для подстановки численных значений в финальную формулу)

\subsubsection{Задача 4.2}

Рассмотрим некоторую задачу классификации. Пусть задано качество 4 моделей a1, a2, a3, a4.
Качество полученных моделей показано в приложении~\ref{app:data-classifiers}. 
Исследователю требуется выбрать наилучшую модель. Для выбора лучшей модели исследовать
требуется попарно сравнить среднее значение качества всех моделей. Может ли исследователь утверждать что какая-то из моделей лучше другой?
Требуется:
\begin{enumerate}
    \item записать задачу формально;
    \item предложить статистику для решения данной задачи;
    \item записать нулевое распределение данной статистики;
    \item записать явно правило принятия решения на основе статистики и нулевого распределения для обеспечения уровня значимости $\alpha=0{,}05$;
    \item проверить гипотезу по записанному критерию, для данных из условия. Противоречат ли они гипотезе?
\end{enumerate}
Все выкладки должны быть сделаны аналитически, без использования компьютера. (допускается использование компьютера для подстановки численных значений в финальную формулу)

\subsubsection{Задача 4.3}

Правительство города М. испытывает систему обнаружения нежелаемых лиц по камерам в метро. В качестве демонстрации работоспособности системы была поставлена цель: найти и задержать опасного рецидивиста по имени Николай Вальный, а также его соучастников. Была собрана выборка из 5000 снимков лица, для которых была проверена гипотеза о несовпадении этого снимка с лицами участников команды Н. Вального. Для 100 фотографий нулевая гипотеза была отвергнута на уровне значимости $\alpha=0{,}05$.
Требуется:
\begin{enumerate}
    \item Определить, в чём недостаток данного подхода и как можно его улучшить.
    \item Предложить наилучший, на ваш взгляд, способ для повышения качества данного решения.
    \item Какую меру качества контролирует данный способ? Какие гарантии он предоставляет?
    \item В чём недостатки данного способа?
    \item Как изменилась мощность при использовании предложенного способа относительно изначальной процедуры проверки?
    \item Известно, что все 5000~фотографий были сделаны для разных людей, и правительство хочет, чтобы система ни в коем случае не упустила преступников. Ответьте на те же вопросы из пунктов 2--4.
    \item Как изменилась мощность при использовании предложенного способа относительно предыдущего способа?
\end{enumerate}
Все выкладки должны быть сделаны аналитически, без использования компьютера (допускается использование компьютера для подстановки численных значений в финальную формулу).

\subsection{Задание 5}


\subsubsection{Задача 5.1}

Даны три объекта \texttt{X\_train}, \texttt{Y\_train}, \texttt{X\_test}\footnote{Архив набора: \PSADPath{labs/lab1/data/regression.zip}.}
\texttt{Y\_train} --- изображение (пиксель кодируется черно-белой компонентой), \texttt{X\_train} --- признаки, соответствующие этому изображению (элемент \texttt{X[i,j]} соответствует набору признаков для пикселя \texttt{Y[i,j]}).
Требуется:
\begin{enumerate}
    \item Провести отбор наиболее значимых признаков и построить регрессию \texttt{X}$\to$\texttt{Y}
    \item Проинтерпетировать признаки (каждый признак является функцией, возможно нелинейной, от значения пикселя)
    \item Получить изображение по \texttt{X\_test} (оцениваться будет качество полученного изображения; ожидается $R^2>0{,}85$ на \texttt{X\_train}, \texttt{Y\_train}).
\end{enumerate}

\subsubsection{Задача 5.2}

Дан архив записей голоса с разметкой акцента.\footnote{\url{https://www.kaggle.com/rtatman/speech-accent-archive}}
Требуется: 
\begin{enumerate}
    \item Отобрать записи, соответствующие странам с минимум 30 респонеднтами в выборке
    \item Получить число пересечений нулевого уровня (zero-crossing) по каждой из записей
    \item Провести ANOVA-анализ по аттрибутам родного языка, пола и возраста для уровня значимости 0.15. Дискретность признака zero-crossing игнорировать.
\end{enumerate}


\subsection{Задание 6}


\subsubsection{Задача 6.1}

Рассмотрим задачу предсказания числа заболевших от экологических анализов (данные --- приложение~\ref{app:data-sick}). Гарантируется, что предсказание описывается линейной моделью.
Так как проведение анализов не является бесплатным, то стоит вопрос о том какие из анализов являются лишними (на уровне значимости $\alpha=0{,}05$) для предсказания линейной модели.
Требуется:
\begin{enumerate}
    \item Записать задачу формально;
    \item Провести отбор признаков линейной модели.
\end{enumerate}
Все выкладки должны быть сделаны аналитически, без использования компьютера. (допускается использование компьютера для подстановки численных значений в финальную формулу)

\subsubsection{Задача 6.2}

В городе Н. правительство решило начать борьбу с превышениями скорости автомобилей. Для выбора стратегии борьбы оно сначала решило провести исследования касательно того, влияет ли используемый водителем автомобиль на среднюю скорость передвижения.
Для этого было сформировано 3 выборки по 20 человек, в каждой из которой людям выдали одинаковые автомобили марок Mitsubishi, Audi и BMW, соответственно. В течение месяца замерялась средняя скорость каждого из автомобилей (данные --- приложение~\ref{app:data-gov}). 
Каждая из пар групп была проверена двувыборочным критерием на равенство распределений, также была проведена поправка на множественность гипотез.
Требуется:
\begin{enumerate}
    \item Описать, в чём недостаток подхода правительства.
    \item Предложить метод для более корректного решения задачи.
    \item Записать формальное условие задачи.
\end{enumerate}
Все выкладки должны быть сделаны аналитически, без использования компьютера. (допускается использование компьютера для подстановки численных значений в финальную формулу)
\section{Лабораторная работа~2}
\label{sec:lab2-tasks}

Второй блок охватывает причинные графы и скрытые марковские модели (гл.~\ref{ch:causality}), временные ряды (гл.~\ref{ch:timeseries}).
Теоретические задачи с полным аналитическим выводом --- по правилам \S\ref{sec:report-requirements}.

\subsection{Задание 1}


\subsubsection{Задача 1.1}

Заданы пары изображений\footnote{\url{https://drive.google.com/file/d/156K2MIDP3UqRc5qkH0VpI9IXhWTA-o6G/view?usp=sharing}}: каждая пара состоит из оригинального (\texttt{original\_id.bmp}) изображения и искажённого (\texttt{modified\_id\_phrase.bmp}).
Также задана пара контрольных изображений \texttt{original\_test}, \texttt{modified\_test}.
В 99% случаев искажение заключается в добавлении белого шума. В 1% случаев искажение заключается в добавлении к изображению скрытого сообщения. 
Алгоритм заключается в следующем:
\begin{enumerate}
    \item У исходной фразы берутся порядковые номера всех символов в порядке английского алфавита (abcz -> 0,1,2,25).
    \item Полученный вектор домножается на неизвестный коэффициент $\alpha$ и складывается с вектором картинки (image = image.flatten() + $\alpha$*v + шум).
    \item Если фраза слишком короткая, искажение продолжается периодически.
\end{enumerate}
Требуется раскодировать фразу из контрольной пары.
\begin{remark}
Предполагается, что искажённые изображения без шума будут найдены с помощью статистических моделей, а не полным перебором. Файлы BMP читаются в среде курса стандартными средствами работы с изображениями.
\end{remark}

\subsubsection{Задача 1.2}

Дана статистика по продаже и отмене бронирований (приложение~\ref{app:data-flights}).\footnote{Полный набор ($n=119386$): \PSADPath{labs/lab2/data/flight_tickets.csv}.}
\begin{enumerate}
    \item Построить граф методом inductive search (допускается использование любого инструмента, в том числе DirectLiNGAM)
    \item Посчитать ACE/ATE между \texttt{previous\_cancellations} (отменой предыдущих заказов) и \texttt{known\_client} (клиент известен системе)
    \item Посчитать ACE/ATE между \texttt{previous\_cancellations} (отменой предыдущих заказов) и \texttt{known\_client} (клиент известен системе) напрямую, без использования графа
    \item Сравнить и проинтерпретировать результат
    \item Предложить свой граф зависимостей, исходя из здравого смысла
    \item Посчитать ACE/ATE между всеми наборами признаков без использования графа и с использованием своего графа. Объяснить разницу.
    \item Допускается использование подвыборок для построения графа.
\end{enumerate}

\subsubsection{Задача 1.3}

Проанализировать консервативность z-критерия для корреляции Пирсона в зависимости от:
\begin{itemize}
    \item Мощности выборки
    \item Проверяемого значения коэффициента корреляции
\end{itemize}
Z-тест для корреляции Пирсона (гл.~\ref{ch:hypothesis}) позволяет проверять не только нулевую корреляцию, но и сопоставление с произвольным значением коэффициента.

\subsection{Задание 2}


\subsubsection{Задача 2.1}

Рассматривается задача тестирования вакцины (данные --- приложение~\ref{app:data-vaccine}) от некоторого вируса. Производство вакцины очень дорогое и затратное по времени, поэтому в день может быть произведена только одна ампула.
Требуется проверить, что вакцина помогает (вероятность заразиться меньше у человека с вакциной чем у человека без вакцины).
Эксперимент ставится следующим образом: каждый день есть два идентичных по здоровью человека. Один из людей принимает вакцину, а второй нет, после чего обоих ставят в одну среду с вирусом. В конце проверяют, кто заразился (в таблице $s$ --- sick, $h$ --- healthy).
Весь мир ждет вакцину от данного вируса, поэтому к руководству института постоянно приходят запросы о сроках завершения тестирования образца. Руководство поручило Вам оценить среднее время, которое понадобится на тестирования данной вакцины. А также провести анализ полученных данных на уровне значимости $\alpha=0{,}05$ и при ошибке второго рода $\beta=0{,}2$.
Требуется:
\begin{enumerate}
    \item Записать задачу формально;
    \item Выполнить оценку среднего количества дней для принятия решения (учесть, что истинные вероятности заразиться с вакциной и без равны $p_1=0{,}2$, $p_2=0{,}5$ соответственно)
    \item Выполнить анализ данных и выяснить работает ли вакцина или нет.
\end{enumerate}
Все выкладки должны быть сделаны аналитически, без использования компьютера.

\subsubsection{Задача 2.2}

Требуется построить ARMA-модель (данные --- приложение~\ref{app:data-arma}).
\begin{enumerate}
    \item Записать задачу формально;
    \item Выписать все формулы аналитически;
    \item Провести вычисления всех параметров модели аналитически.
\end{enumerate}
\begin{remark}
Порядки $p$ и $q$ модели ARMA допускается оценить при помощи компьютера. Остальные выкладки --- аналитически, с опорой на гл.~\ref{ch:timeseries}.
\end{remark}
Все выкладки должны быть сделаны аналитически, без использования компьютера.

\subsection{Задание 3}


\subsubsection{Задача 3.1}

Дан корпус Nerus --- предложения на русском языке с разметкой частей речи.\footnote{\url{https://github.com/natasha/nerus}}
Требуется:
\begin{enumerate}
    \item Рассмотреть последовательность частей речи как марковскую модель. Определить оптимальный порядок марковской модели.
    \item Обучить скрытую марковскую модель по выборке. Оценить точность предсказания частей речи, посчитать энтропию на выборке.
\end{enumerate}
\begin{remark}
В целях ускорения эксперимента рекомендуется взять первые 10~МБ текста из корпуса.
\end{remark}

\subsubsection{Задача 3.2}

Дан корпус Nerus --- предложения с разметкой именованных сущностей.\footnote{\url{https://github.com/natasha/nerus}} 
Требуется:
\begin{enumerate}
    \item Рассмотреть последовательность именованных и неименованных сущностей как марковскую модель. Определить оптимальный порядок марковской модели.
    \item Обучить скрытую марковскую модель по выборке. Оценить точность предсказания, является ли слово именованной сущностью или нет, посчитать энтропию на выборке.
\end{enumerate}
\begin{remark}
В целях ускорения эксперимента рекомендуется взять первые 10~МБ текста из корпуса.
\end{remark}

\subsection{Задание 4}


\subsubsection{Задача 4.1}

На основе следующей языковой модели требуется:
\begin{enumerate}
    \item Выбрать наиболее вероятную последовательность слов, которая начинается со слова \texttt{<start>} и заканчивается словом \texttt{<end>}.
    \item Выбрать последовательность токенов, которая начинается со слова \texttt{<start>} и заканчивается словом \texttt{<end>} и имеет максимальный score $\frac{1}{l}\sum_{i=1}^{l}\log p(w_i\mid w_{i-1})$, где $l$ --- длина полученной последовательности.
\end{enumerate}
\begin{verbatim}
p(обучение|<start>)=0.14
p(<end>|машинный)=0.5
p(<end>|обучение)=0.5
p(не|обучение)=0.5
p(без|возможно)=0.5
p(без|<start>)=0.14
p(не|<start>)=0.14
p(<end>|без)=0.5
p(<end>|не)=0.5
p(обучение|машинный)=0.5
p(<end>|статистика)=1.0
p(возможно|не)=0.5
p(возможно|<start>)=0.14
p(машинный|<start>)=0.29
p(статистика|<start>)=0.14
p(статистика|без)=0.5
p(<end>|возможно)=0.5
\end{verbatim}
Все выкладки должны быть сделаны аналитически, без использования компьютера.
\begin{remark}
Решение полным перебором не оценивается в полный балл.
\end{remark}

\subsubsection{Задача 4.2}

Получить аналитическую формулу для оценки правдоподобия последовательности слов из словаря мощности $d$ на основе языковой модели, в случаях:
\begin{enumerate}
    \item Если все слова независимы и распределены равномерно.
    \item Если известно распределение $p(w_i\mid w_{i-1})=\mathrm{Dir}(\mu(w_{i-1}))$, где $\mathrm{Dir}$ --- распределение Дирихле с вектором параметров $\mu(w_{i-1})$ размерности $d$. Компоненты $\mu_k=\exp(-|k-j|)$, где $j$ --- номер слова $w_{i-1}$ в словаре, $k$ --- индекс слова в словаре.
\end{enumerate}
Все выкладки должны быть сделаны аналитически, без использования компьютера.
Подсказка по п.~2: удобно работать с логарифмом правдоподобия. Зависимостями между словами на расстоянии больше двух пренебречь. В one-hot векторах в распределении Дирихле предположить, что вероятность слова $i$ равна $1-\varepsilon$, а всех остальных --- $\varepsilon/|d|$.
\section{Архив контрольных работ}
\label{sec:exam-tasks}

Типовые формулировки из печатных бланков 2023--2026 годов. На экзамене вариант определяется по формуле из бланка. Численные значения и p-value теста Шапиро--Уилка могут отличаться между частями I и II.

\subsection{Сравнение средств для мытья окон}

\subsubsection{Парные измерения в разных квартирах}

При тестировании нового средства для мытья окон были собраны показания коэффициента проникновения света в помещениях после мытья стекол одного из небоскребов Москвы. Для исследования было выбрано 10~квартир в небоскребе. Тестирования проводились в одних и тех же квартирах в разные месяцы. После мытья для каждого окна был измерен коэффициент интенсивности проходящего света (чем больше, тем лучше). Верно ли, что новое средство более эффективно?

При проверке гипотезы можно воспользоваться p-value теста Шапиро--Уилка из бланка. Пример пар измерений (старое, новое) для 10~квартир --- $(10.3, 10.8)$, $(10.1, 9.2)$, $(10.1, 9.1)$, $(9.7, 10.0)$, $(10.1, 9.3)$, $(9.8, 9.1)$, $(10.9, 7.6)$, $(9.9, 9.5)$, $(9.3, 9.1)$, $(9.1, 10.1)$.

\subsubsection{Парные измерения на половинах стекла}

При тестировании нового средства для мытья окон была собрана выборка из стекол одного из небоскребов Москвы. Каждое стекло было условно разделено на две половины. Первая половина была помыта при помощи нового средства, вторая --- старым. После мытья для каждой половины был измерен коэффициент интенсивности (или прозрачности) проходящего света. Верно ли, что новое средство более эффективно? Данные --- те же 10~пар, что в предыдущей задаче.

\subsubsection{Разности коэффициента проникновения}

При тестировании нового средства для мытья окон были собраны показания разности коэффициента проникновения света в помещениях до и после мытья стекол одного из небоскребов Москвы. Для исследования было выбрано 10~квартир. Пример разностей для старого средства --- 10.3, 10.1, 10.1, 9.7, 10.1, 9.8, 10.9, 9.9, 9.3, 9.1. Для нового --- 10.8, 9.2, 9.1, 10.0, 9.3, 9.1, 7.6, 9.5, 9.1, 10.1.

\subsection{Контроль массы грузиков}

На производстве грузиков для весов максимально допустимое отклонение составляет 1\% от номинала 10~г. При ревизии из партии отобрано $n=10$ (или $n=100$) грузиков. Можно ли сказать, что партия удовлетворяет требованиям? Пример весов при $n=10$ --- 9.2, 8.9, 9.8, 9.1, 9.1, 10.1, 9.9, 9.9, 9.5, 9.4.

\subsection{Сравнение моделей детекции}

Проводится тестирование новой модели детекции машинной генерации. Для оценки качества проводится сравнение моделей на 10~выборках (чем выше значение качества, тем лучше). Верно ли, что новая модель качественнее старой? Пример пар (старая, новая) --- $(90.3, 90.8)$, $(89.1, 88.2)$, $(90.1, 89.1)$, $(88.7, 90.0)$, $(90.1, 89.3)$, $(90.8, 90.1)$, $(90.9, 7.6)$, $(87.9, 89.5)$, $(86.3, 90.1)$, $(86.1, 90.1)$.

\subsection{Препараты и артериальное давление}

При тестировании трёх новых средств от болезни~X было сформировано три группы людей по 5~человек, которые принимали препараты $B$, $C$, $D$. Отдельная группа проходила классическое лечение~$A$. Показатель --- артериальное давление (чем ниже, тем лучше). Верно ли, что новые средства эффективнее классического лечения? Предложите два наиболее мощных критерия с контролем ошибок.

\subsection{Согласованность моделей}

\subsubsection{Асессорские ранги}

Задана ранжирующая модель ML~--- $m_1$. Проведено дообучение $m_1$ в модель $m_2$. Оценка качества --- асессорские ранги от 0 до 100. Требуется проверить согласованность ответов асессоров. Поставить задачу формально, указать статистику, критерий и нулевое распределение.

\subsubsection{ROC AUC}

Задана модель ML для бинарной классификации~--- $m_1$, дообучена в $m_2$. Оценка качества --- метрика ROC AUC от 0 до 1. Требуется проверить согласованность моделей $m_1$ и $m_2$. Поставить задачу формально, указать статистику, критерий и нулевое распределение.

\subsection{Помощник покупателя}

Рассматривается задача оценки недополученной прибыли компании. Каждый объект --- сделка (успешная или неуспешная). В компании внедряют помощника покупателя. Требуется проверить, что при внедрении помощника доля успешных сделок станет выше. Поставьте задачу формально и оцените необходимый объём сделок для оценки значимости помощника при $\alpha=0{,}05$, $\beta=0{,}2$. Любые гипотезы о распределении данных должны быть выписаны в решении.

\subsection{Логистические модели}

\subsubsection{Поломка электропоезда}

Пусть задана модель предсказания поломки электропоезда $g(x)=\sum_i x_i w_i$, где третий признак описывает давность техосмотра ($0$ --- менее месяца назад, $1$ --- более месяца). Известно $w_3=2$. Определите увеличение шансов поломки для поездов с просроченным техосмотром.

\subsubsection{Важность признака техосмотра}

При той же модели поломки поезда, при $\widehat{D}_{w_3}=4$, определите важность признака давности техосмотра.

\subsubsection{Заболевание и лейкоциты}

Пусть задана модель предсказания заболевания $g(x)=\sum_i x_i w_i$, где третий признак --- уровень лейкоцитов от 0 до 20. Известно $w_3=2$. Определите увеличение шансов заболевания при увеличении уровня лейкоцитов на~1.

\subsubsection{Доверительный интервал для коэффициента}

При модели заболевания, при $w_3=2$ и $\widehat{D}_{w_3}=4$, определите 95\%-й доверительный интервал для $w_3$.

\subsection{Временные ряды}

\subsubsection{STL-разложение}

Выполнить STL-разложение для ряда. Проверить ряд на стационарность и все причины стационарности или нестационарности.

\subsubsection{Подбор ARMA}

Выполнить настройку параметров ARMA-модели для ряда. Проверить ряд на стационарность и все причины стационарности или нестационарности.

\subsection{Причинные графы}

В заданном графе причинности найти d-разделимое множество (или все d-разделимые множества) для пары $(B, G)$.

\subsection{Скрытые марковские модели}

\subsubsection{Самочувствие и диагноз}

Доктор опрашивает людей о самочувствии (normal, dizzy, cold). Скрытые состояния --- healthy или fever. Требуется восстановить вероятности переходов по таблице наблюдений.

\begin{center}
\small
\begin{tabular}{lcccccccccc}
Самочувствие & N & D & D & C & C & D & N & C & N & N \\
Диагноз & F & F & H & F & H & F & F & H & F & H \\
\end{tabular}
\end{center}

Обозначения --- H healthy, F fever, N normal, D dizzy, C cold.

\subsubsection{Гласные и согласные}

Задана модель генерации английского текста с скрытыми состояниями vowel и conson. Требуется восстановить вероятности переходов по тексту (фрагмент из бланка контрольной работы).

\section{Требования к отчёту}
\label{sec:report-requirements}

Ожидается отдельный ноутбук на задачу, формулировка $H_0$ или модели, диагностика до основного теста, p-value или интервал и размер эффекта, поправка при множественности, содержательный вывод в 3--5 предложениях.
Теоретические задачи требуют полного аналитического вывода: все переходы описаны, выкладки и преобразования явно посчитаны (допускается подстановка численных значений в финальные формулы).
Обоснованность выбора критериев, учёт поправки на множественности, предобработка реальных данных и учёт погрешности симуляций --- обязательные критерии оценки.

\section{Типичные ошибки}

Смешение доверительного и предсказательного интервала. Попарные $t$-тесты вместо дисперсионного анализа. Интерпретация скорректированного $p$ как вероятности $H_0$. Пошаговый отбор с классическими $p$ на той же выборке. $\chi^2$ при малых ожидаемых частотах без критерия Фишера.
